\documentclass[11pt,letterpaper]{article}
\usepackage[margin=0.87in,top=0.78in,bottom=0.85in]{geometry}
\usepackage{amsmath,amssymb,booktabs}
\usepackage[hidelinks]{hyperref}
\hypersetup{pdftitle={Class-group exponent seven in imaginary quadratic fields and orders},
  pdfauthor={nunavipahe60-star},pdfsubject={A classification under Dirichlet GRH}}
\pdfinfoomitdate=1
\pdfsuppressptexinfo=15
\pdftrailerid{}
\setlength{\parindent}{0pt}
\setlength{\parskip}{4.5pt plus 0.4pt minus 0.4pt}
\setlength{\emergencystretch}{1.5em}
\newcommand{\file}[1]{\texttt{\detokenize{#1}}}
\newcommand{\Cl}{\operatorname{Cl}}
\newcommand{\Pic}{\operatorname{Pic}}
\newcommand{\F}{\mathbb F}
\newcommand{\Z}{\mathbb Z}
\newcommand{\sect}[1]{\section*{#1}}
\newcommand{\arxivref}[1]{\href{https://arxiv.org/abs/#1}{arXiv:#1}}
\begin{document}
\begin{center}
{\Large\bfseries Class-group exponent seven in imaginary quadratic\\[2pt]
fields and orders}\\[8pt]
{\large nunavipahe60-star}\\[3pt]
26 November 2026
\end{center}

\textbf{Abstract.}
Assuming Dirichlet GRH, we determine all imaginary quadratic fields and orders
whose class group has exponent seven. An explicit formula for the odd quadratic
character bounds the least split prime by \(4597\). A finite trace sieve then
excludes the entire outstanding discriminant interval. A second computation
eliminates the discriminant between two norm equations and exhausts the
resulting divisors on a CPU.

\sect{1. Classification and reduction to prime discriminants}
\textbf{Theorem.} Under Dirichlet GRH, exactly \(33\) imaginary quadratic fields
have ideal class group of exponent seven. Exactly \(40\) imaginary quadratic
orders, including those \(33\) maximal orders, have proper Picard group of
exponent seven. The largest absolute discriminant is \(118843\). The two
discriminants \(-63499,-118843\) give \(C_7^2\), and the other \(38\) give \(C_7\).

Elsenhans, Kl\"uners and Nicolae (EKN) classified the fields through
\(L=3.1\cdot10^{20}\) and gave the conditional upper bound
\(Q_0=3.9\cdot10^{30}\) \cite[Theorems 1, 3 and Tables 1--2]{EKN}.
Aiken, Ross, Shvydkoy and Xue (ARSX) reduce completeness of their Table~A.7
to excluding fundamental discriminants in this interval
\cite[Theorem 1.3(6)]{ARSX}. We use Dirichlet GRH throughout. In particular,
the quadratic-field zeta functions needed at the upper endpoint factor as
\(\zeta(s)L(s,\chi_D)\).

Suppose a field of exponent seven with fundamental discriminant \(D\) lies
in the gap. Its class group is \(C_7^r\), so its \(2\)-rank is zero. By genus theory
\(D\) has one prime-discriminant factor. The even possibilities \(-4,-8\)
lie below \(L\). Hence \(D=-q\) with \(q\equiv3\bmod4\) prime and \(L<q\le Q_0\).
If \(p\) splits, a prime ideal \(\mathfrak p\mid p\) satisfies
\(\mathfrak p^7=((x+y\sqrt{-q})/2)\). Its norm gives
\begin{equation}
 4p^7=x^2+qy^2,\qquad x,y\ge1,\qquad x\equiv y\pmod2 .
 \label{eq:norm}
\end{equation}
Indeed \(y=0\) would make \(p^7\) a square, and \(x=0\) would force
\(q\mid4p^7\). Absolute values give the indicated signs. Thus \(q<4p^7\).
As \(4\cdot691^7<L<4\cdot701^7\), the least split prime \(p_1\) is at least
\(701\). In particular 2 is inert, so \(-q\equiv5\bmod8\) and \(q\equiv3\bmod8\).
At \(p_1=701\), \(4p_1^7<4L\) forces \(y=1\). The sieve in Section~4 includes
this endpoint; its separate integer check already terminates at the prime \(137\).

\textbf{Corollary for nonmaximal orders.}
Extension of ideals maps the Picard group at conductor \(f\) onto the field
class group. ARSX's conductor enumeration includes fields of exponent one
and seven \cite[\S\S5--6]{ARSX}. Once the fundamental gap is excluded, its
seven nonmaximal discriminants are
\(-284,-343,-604,-892,-1372,-1852,-1948\).
Together with the separate exponent-six result for Theorem~1.3(5) \cite{six},
this completes Tables~A.1--A.8 for \(1\le E\le8\) under Dirichlet GRH.
The exponent-seven proof uses only its own discriminant and prime bounds.

\newpage
\sect{2. A Fourier-positive certificate with wide support}
Give \(g\) the signed integer values of a certificate on \(m\) equal cells
of \([0,\log T)\), extending by zero. Put
\[
 f(u)=\frac{\int_{\mathbb R} g(t)g(t+u)\,dt}{\int_{\mathbb R}g(t)^2\,dt},
 \qquad h(z)=\int_{\mathbb R}f(u)e^{izu}\,du.
\]
The even, piecewise linear \(f\) has \(f(0)=1\), \(|f|\le1\) and support
\([-\log T,\log T]\). On the real axis \(h\) is a normalized squared
Fourier modulus. Its entire extension decays as
\(O_T((1+|\Re z|)^{-2})\) on bounded horizontal strips. For the primitive
odd character \(\chi=\chi_{-q}\), the GRH zero sum is therefore nonnegative.
The Guinand--Weil formula \cite[Lemma 9]{CMQR} gives
\begin{equation}
 2\sum_{n<T}\frac{\Lambda(n)\chi(n)f(\log n)}{\sqrt n}
 \le \log(q/\pi)+\psi(3/4)+J_f,\qquad
 J_f=\int_0^\infty\frac{e^{-u/2}(1-f(u))}{\sinh u}\,du .
 \label{eq:gw}
\end{equation}
Here \(\psi=\Gamma'/\Gamma\). Fourier inversion in the integral for
\(\Re\psi(3/4+it/2)-\psi(3/4)\) yields the displayed weight.

Let \(A_p,C_p\) denote the left-hand contribution of \(p\) when
\(\chi(p)=-1,+1\), respectively. Thus \(A_p=2\log p\sum_{p^j<T}
(-1)^j f(j\log p)p^{-j/2}\), while \(C_p\) uses \(+1\) for every power.
Write \(B_p=A_p-C_p\) and \(b^+=\max(0,b)\). Since \(T<q\), every prime
in the sum is unramified. With no split prime through \(S<T\), the left
side is at least \(\sum_{p<T}A_p-\sum_{S<p<T}B_p^+\). For \(q\le Q\), a
positive lower bound for
\[
 \sum_{p<T}A_p-\sum_{S<p<T}B_p^+-\log Q+\log\pi-\psi(3/4)-J_f
\]
therefore proves \(p_1\le S\). The \(256\)-cell certificates below use
\(Q=Q_0\) in the first column and \(Q=4S_{\rm prev}^7\) thereafter.
\begin{center}
\begin{tabular}{@{}lrrrrrr@{}}
\toprule
Step&1&2&3&4&5&6\\
\midrule
Support \(T\)&20550&18893&14972&13616&13164&13073\\
Bound \(S\)&4597&3643&3313&3203&3181&3169\\
Margin, rounded down&.059634&.065825&.022050&.028155&.060095&.030033\\
\bottomrule
\end{tabular}
\end{center}
Thus \(Q_* =4\cdot3169^7=12\,838\,540\,801\,993\,873\,675\,201\,156\) is valid.
To certify \(p_k\) for \(k\ge2\), fix the actual \(a=p_1\) within its earlier
bound. Deduct \(B_a\), the largest \(k-2\) positive losses with \(a<p\le S\)
(or all if fewer), and the tail above \(S\). Use \(q\le\min(Q,4a^7)\) and
take the worst result over admissible \(a\). Selected bounds are shown below.
\begin{center}
\begin{tabular}{@{}lrrrrrrrr@{}}
\toprule
\(k\)&1&2&3&6&16&32&64&128\\
\midrule
All ranks&3169&3181&3187&3209&3301&3491&4943&8513\\
Maximum over Section~3 bands&3089&3119&3121&3169&3257&3449&4639&6581\\
\bottomrule
\end{tabular}
\end{center}
The verifier encloses all prime powers and rational autocorrelation values
by outward-rounded intervals. For the first cell of width \(d=\log(T)/m\),
the integral is at most \(1-f(d)\). On each later cell it uses 32 subcells,
multiplying the decreasing weight at the left endpoint by the exact integral
of \(1-f\). Beyond \(\log T\) it adds \(4/(3T^{3/2}(1-T^{-2}))\).

\newpage
\sect{3. Short Euler sums and the possible ranks}
For the independent two-prime computation we use the class number formula
\(7^r=\sqrt q\,L(1,\chi)/\pi\). We record the finite inequalities from
Lamzouri, Li and Soundararajan (LLS) \cite[\S5]{LLS} to specify both the
truncation and the direction of the rounding. For \(X\ge100\), let
\(t=\log X,\ s=\sqrt X\), and put
\[
 a_\epsilon=\frac{1/t-2\epsilon/(st^2)}{(1+\epsilon/s)^2},\quad
 b_+=\frac{1+1/t}{(s+1)t},\quad
 b_-=-\frac{1+(1+1/s)/t}{(s-1)t}\qquad(\epsilon=\pm1).
\]
Define the weighted prime-power sum
\[
 W_X(a)=\sum_{p^j\le X}\chi(p)^j
 \left[\frac{1-j\log(p)/t}{j p^j}
       +a\log p\bigl(p^{-j}-X^{-1}\bigr)\right].
\]
The LLS upper bound for \(\log L(1,\chi)\) is
\(W_X(a_+)+b_+\log(q/\pi)\), and the lower bound is
\(W_X(a_-)+b_-\log(q/\pi)\). All displayed summand weights are
nonnegative at the cutoffs used here. Even powers have sign \(+1\), since
\(q>X\). For the lower bound give all odd powers sign \(-1\). For the upper
bound retain \(-1\) at the known inert primes and give other odd powers
sign \(+1\). With \(q\le Q_*\), the conductor term is extremized at \(Q_*\)
for each inequality. Taking \(X=2100,4000\), respectively, and allowing
701 as the first possible split prime, outward rounding gives
\[
 0.10919<L(1,\chi)<0.18666,\qquad
 \text{in particular}\quad 0.1084<L(1,\chi)<0.1875 .
\]
Using \(L<q\le Q_*\) gives \(7^{10}<h(-q)<7^{14}\), so \(r\in\{11,12,13\}\).
The initial certificate also permits the looser cap \(4\cdot4339^7\), with
uniform bounds \(0.10840844<L(1,\chi)<0.18752520\).

For each \(r\), a trial lower bound on \(q\) excludes further split primes
by \eqref{eq:norm}. Reevaluate the LLS upper bound for \(h\) with those
inert signs fixed. The lower and upper bounds for \(h\) are increasing in
\(q\), since \(1/2+b_\pm>0\). Testing successive endpoints gives these
closed outer intervals. Applying Section~2 within each gives the prime bounds.
\begin{center}
\begin{tabular}{@{}crrcrr@{}}
\toprule
\(r\)&lower \(q\)&upper \(q\)&range for \(p_1\)&upper \(p_2\)&\((y_{\max},v_{\max})\)\\
\midrule
11&\(1.259\cdot10^{21}\)&\(3.044\cdot10^{21}\)&853--2371&2389&\((36,37)\)\\
12&\(7.873\cdot10^{22}\)&\(1.538\cdot10^{23}\)&1531--2713&2729&\((7,7)\)\\
13&\(4.807\cdot10^{24}\)&\(7.743\cdot10^{24}\)&2767--3089&3119&\((1,1)\)\\
\bottomrule
\end{tabular}
\end{center}
\textbf{Relations among small prime classes.}
In rank one, a trivial small prime class gives \(q\le4p_i\). Otherwise
the second prime class equals the first to a power \(\pm1,\pm2,\pm3\).
Orienting the ideals by conjugation yields a nonrational principal
generator and \(q\le4p_1^3p_2\le404\,939\,719\,793\,716<L\).

For high rank the abstract group alone cannot force a relation of signed
degree at most six among 128 labels. Take \((1,z,z^2,z^3)\) for all
\(z\in\F_{7^m}\), in \(\F_7\oplus\F_{7^m}^3\). In any such relation the
first coordinate forces equal numbers of terms on both sides, at most three.
Their first three power sums agree. Newton's identities recover the two
multisets, proving the relation trivial. This gives 343 labels in
dimension 10 or 2401 in dimension 13. Holmin et al.\ already excluded
ranks 3 through 7 for this prime under GRH \cite[\S1.1.1]{HJKMP}.

\newpage
\sect{4. Finite exclusion and verification of the orders}
\textbf{The single-trace search.}
Let \(p=p_1\) and \(U_p=\lfloor\sqrt{4p^7-L}\rfloor\).
Equation~\eqref{eq:norm} forces \(1\le x\le U_p\). For every odd prime
\(\ell<p\), primality of \(q>\ell\) and the minimality of \(p\) imply
\(\chi(\ell)=-1\). Therefore
\[
 F_p(x)=x^2-4p^7=-qy^2,\qquad
 \left(\frac{F_p(x)}{\ell}\right)\in\{0,-1\}.
\]
A nonzero square rules out the entire trace for every possible \(y\).
Zero remains admissible because \(\ell\) can divide \(y\).

For each \(701\le p\le4597\) prime, compute the admissible residue classes
modulo \(M=3\cdot5\cdot7\cdot11\cdot13\cdot17\cdot19\cdot23=111\,546\,435\).
Every positive trace left by these predicates occurs exactly once as
\[
 x=r+(32j+b)M,\qquad 0\le r<M,\quad j\ge0,\quad 0\le b<32 .
\]
For each remaining odd prime through 251 a 32-bit mask tests all lanes
\(b\) of a tile. Intersect the masks, clip the final tile at \(U_p\),
and send every retained \(x\) to the host. A separately implemented
CUDA C++ kernel compiled with NVRTC uses 29 in place of 23 in its wheel,
constructs residues directly by the Chinese remainder theorem, forms local
tables from sets of squares and reverses the other test primes. The primary
kernel is written in PTX. Both kernels use the same mathematical predicate.

\begin{center}
\begin{tabular}{@{}lrrr@{}}
\toprule
Run&Prime rows&Raw trace positions \(\sum_p U_p\)&Retained through 251\\
\midrule
PTX through 4597&497&\(1\,622\,545\,300\,159\,089\)&5\\
NVRTC through 4597&497&\(1\,622\,545\,300\,159\,089\)&5\\
PTX through 3169&324&\(311\,248\,928\,889\,382\)&0\\
\bottomrule
\end{tabular}
\end{center}
Both wide runs returned the same five pairs. Each row below exhibits a
prime \(\ell<p\) and \(c\) such that \(F_p(x)\equiv c^2\not\equiv0\bmod\ell\).
\begin{center}
\begin{tabular}{@{}rrrrr@{}}
\toprule
\(p\)&\(x\)&\(\ell\)&\(F_p(x)\bmod\ell\)&\(c\)\\
\midrule
3209&\(2\,044\,007\,472\,677\)&263&51&83\\
3853&\(3\,871\,258\,885\,871\)&271&2&96\\
4129&\(5\,310\,890\,329\,408\)&257&239&53\\
4133&\(7\,753\,344\,152\,825\)&269&24&37\\
4561&\(3\,028\,955\,847\,366\)&257&42&82\\
\bottomrule
\end{tabular}
\end{center}
This exhausts the range supplied by the first Fourier row and proves the
field assertion. The six-row analytic bound alternatively requires just
the short run, which took \(63.84\) seconds on a GeForce RTX 2060.

\textbf{Positive entries.}
For every discriminant in Table~A.7 we enumerated primitive reduced positive
binary quadratic forms, multiplied the associated ideal lattices to the
seventh power, and tested that the reduced norm form of each power has
minimum one. All \(38\cdot7+2\cdot49=364\) ideal classes passed. The
class numbers and the seven conductor discriminants of Section~1 then give
the stated group structures and order count.

\newpage
\textbf{An independent CPU exclusion by two split primes.}
Let \(p<s\) be the first two split primes, with data \((x,y)\) and \((u,v)\)
in \eqref{eq:norm}. Write \(g=\gcd(y,v),\ y_0=y/g,\ v_0=v/g\).
Eliminating \(q\) gives
\[
 (v_0x-y_0u)(v_0x+y_0u)=4V,\qquad
 V=v_0^2p^7-y_0^2s^7\ne0 .
\]
The inequality follows because \((p/s)^7\) has odd valuation at \(p\).
Both factors are even by the trace parities. Thus a necessary solution
gives positive divisors \(a=|v_0x-y_0u|/2,\ b=(v_0x+y_0u)/2\) of
\(K=|V|\) with \(ab=K,\ a\le b\). Conversely the candidate traces are
recovered by exact division from
\[
 v_0x=b+\operatorname{sgn}(V)a,\qquad
 y_0u=b-\operatorname{sgn}(V)a .
\]
For each \((p,s,y,v)\) use the rank band to bound \(x,u\) by integer
square roots. The induced bounds on \(v_0x+y_0u\) restrict
\(a=2K/(v_0x+y_0u)\). Check \(q\bmod8\) and the compatibility of the
two equations at \(3,5,7,11,13\). For every remaining tuple, factor \(K\)
using FLINT, prove each prime factor with \file{is_prime}, multiply back,
and enumerate every divisor in the interval. Retain integral
representations with in-band \(q\equiv3\bmod8\). Finally require
\(({-q}/\ell)=-1\) for every odd prime \(\ell<p\).
\begin{center}
\begin{tabular}{@{}rrrrr@{}}
\toprule
Rank&Locally admissible tuples&Factorizations&Representations&Survivors\\
\midrule
11&542464&493357&13115&0\\
12&44026&39339&5802&0\\
13&860&743&655&0\\
\midrule
Total&587350&533439&19572&0\\
\bottomrule
\end{tabular}
\end{center}
The largest \(K\) has 89 bits. The full CPU run took \(877.91\) seconds
on a six-core Core i7-10750H. As a positive control the same search
recovers \(q=5923\) from \(p=19,s=29\), with
\((x,y)=(13128,758)\) and \((u,v)=(2243,3413)\).
This arithmetic check uses the analytic bounds of Sections~2--3.
The source, interval witnesses, per-prime records and reproduction commands
accompany the paper.

\textbf{Acknowledgment.}
The project was coordinated by nunavipahe60-star in collaboration with
Claude, Gemini and Codex. The claims ledger records the scope of each review.

\begingroup
\small
\begin{thebibliography}{9}
\setlength{\itemsep}{1pt}
\bibitem{EKN} A.-S. Elsenhans, J. Kl\"uners and F. Nicolae,
Imaginary quadratic number fields with class groups of small exponent,
\emph{Acta Arith.} \textbf{193} (2020), 217--233.
\href{https://doi.org/10.4064/aa180220-20-3}{doi:10.4064/aa180220-20-3}.
\bibitem{ARSX} D. Aiken, E. Ross, D. Shvydkoy and H. Xue,
Boundary CM points and class groups of small exponent, 2026.
\arxivref{2605.26020}.
\bibitem{CMQR} E. Carneiro, M. B. Milinovich, E. Quesada-Herrera and
A. P. Ramos, Fourier optimization, the least quadratic non-residue, and the
least prime in an arithmetic progression, 2024, revised 2025.
\arxivref{2404.08380v2}.
\bibitem{LLS} Y. Lamzouri, X. Li and K. Soundararajan, Conditional bounds
for the least quadratic non-residue and related problems, corrected version,
2016. \arxivref{1309.3595v3}.
\bibitem{HJKMP} S. Holmin, N. Jones, P. Kurlberg, C. McLeman and K. L. Petersen,
Missing class groups and class number statistics for imaginary quadratic
fields, 2015. \arxivref{1510.04387}.
\bibitem{six} nunavipahe60-star, \emph{Class-group exponent six in imaginary
quadratic fields and orders}, companion repository, Zenodo, 2026.
\href{https://doi.org/10.5281/zenodo.24489676}{doi:10.5281/zenodo.24489676}.
\end{thebibliography}
\endgroup
\end{document}
